\subsection{THETA3-C38. A terminating finite certificate for simultaneous theta differences}
\label{C38}

The compactness proof of C37 admits a finite verification procedure.
The conclusion here is an algorithm with a proof of termination for
computable fixed theta data satisfying C37, rather than a universal
closed-form numerical bound.  It retains the same-prime graph
measures, arbitrary unbounded representatives, and all phase
parameters.  In particular it does not replace these graph measures
by product Haar measure.

\paragraph{Computable input and a scalar positivity test.}
Retain $T=V/\Lambda$, $F$, $f$, and the finite sample families of
C37.  Fix real coordinates $T=\R^d/\Z^d$, $d=2\dim_{\C}V$,
and put $r=f-1$.  Assume the period basis, $H$, and the fixed
theta function are computable with certified uniform error and
derivative bounds on compact sets.  Gaussian theta-series tail
estimates, such as the estimates already retained in Q2 and A12,
provide such data.  The finite torsion labels, their complex
coordinates, and reductions of fixed trace-dual elements must also
be computable.  An unbounded computably enumerable family of
allowed primes is the final input.  All these are hypotheses about
specified data; this statement does not supply a missing fixed
baseline or an unspecified archive prime family.

Choose computable $t_2,\ldots,t_f$ satisfying (C37d).  For example,
if $e_1$ is one vector of the real period basis, take
$t_i=a_i\sqrt{p_i}\,e_1$, with nonzero rational $a_i$ and distinct
rational primes $p_i$.  A nonzero signed sum of these coefficients
is not rational: the independent sign automorphisms of their
multiquadratic field prove this immediately.  Thus distinct subset
sums differ outside $\Lambda\otimes\Q$.  These choices may be
reduced modulo $\Lambda$ or made small by the rational $a_i$.

Let $\Sigma$ be the compact parameter space consisting of unit
tangent vectors $u$ and independent unit scalars $c_J$ for
$J\subset\{2,\ldots,f\}$.  For $\theta=(u,(c_J))\in\Sigma$
write, on $\mathcal T=T\times(S^1)^r$,
\[
 A_\theta(y,z)=\sum_J c_Jz^Jd(U_{t_J}U_yF)_0(u),
 \qquad
 g_\theta(y,z)=
       \bigl(|A_\theta(y,z)|^2-|A_\theta(0,1)|^2\bigr)^2.
 \tag{C38a}
\]
The function $A_\theta$ is a section, but $g_\theta$ is a
well-defined smooth nonnegative scalar function on $\mathcal T$.
Its coefficients and uniform bounds for its derivatives, including
derivatives in the parameters, can be computed from the stated
theta data.  Let $M_0$ be a certified uniform bound for
$|A_\theta|$ and $|A_\theta(0,1)|$.  One may use the conservative
bound $B_g=4M_0^4$ for $\|g_\theta\|_\infty$.

For every closed subgroup $G\subset\mathcal T$ projecting onto
$T$ and every $\theta\in\Sigma$, C37 proves
\[
                     \int_G g_\theta\,dm_G>0.
 \tag{C38b}
\]
Indeed a zero integral would make the norm of $A_\theta$ constant
on $G$, the exact impossibility proved there.

These integrals have a common positive lower bound, even as $G$
varies.  A direct compactness proof is useful.  A Haar subgroup
measure has Fourier coefficients $b(\xi)\in\{0,1\}$, equal to
one precisely on a subgroup of $\Z^{d+r}$.  Projection onto $T$
is equivalent to
\[
       b(\gamma,0)=0\quad(0\ne\gamma\in\Z^d).
 \tag{C38c}
\]
The set of these patterns is closed in
$\{0,1\}^{\Z^{d+r}}$: the subgroup conditions and (C38c) are
closed finite-coordinate conditions.  It is therefore compact,
as can also be proved by a diagonal subsequence.  Convergence of
patterns gives weak convergence of the associated Haar measures,
first on trigonometric polynomials and then on continuous
functions.  Joint continuity with the compact parameter space
$\Sigma$ and (C38b) give a number $\Delta>0$ with
\[
                \int_Gg_\theta\,dm_G\ge\Delta
                       \quad\text{for every }G,\theta.
 \tag{C38d}
\]
The algorithm below discovers a positive lower bound without
being given $\Delta$.

\paragraph{Exactly which finite relation patterns are possible.}
Let $W\subset\Z^{d+r}$ be a finite symmetric set containing zero.
For a proposed pattern $b:W\to\{0,1\}$ with $b(0)=1$, form
the integer subgroup
\[
              D_b=\sum_{\xi\in W,\ b(\xi)=1}\Z\xi.
 \tag{C38e}
\]
Do not saturate this subgroup.  The pattern is the restriction of
a Haar subgroup measure projecting onto $T$ if and only if
\[
 \begin{cases}
 \xi\notin D_b&\text{for every }\xi\in W\text{ marked }0,\\
 D_b\cap(\Z^d\times\{0\})=\{0\}.&
 \end{cases}
 \tag{C38f}
\]
Necessity follows because $D_b$ is contained in the annihilator
of the subgroup.  For sufficiency take $G_b$ to be the
annihilator of $D_b$.  The double-annihilator identity for
subgroups of $\Z^{d+r}$, proved for example by integer normal
forms, says that its Haar pattern is exactly membership in
$D_b$.  Thus it has the prescribed values on $W$.  The second
line of (C38f) says that its projection to $T$ has no nontrivial
annihilating character, and hence is all of $T$.  This proves
both directions.  Nonsaturated $D_b$ are necessary, since $G_b$
can be disconnected.  Membership, intersection, and witnesses for
failure of the second line are finite integer linear-algebra
computations using Hermite or Smith normal forms.

Two small examples explain the qualifications.  In
$\Z^2=\Z_{\rm base}\oplus\Z_{\rm circle}$, the subgroup
$\Z(0,2)$ is admissible and gives
$G=T\times\{1,-1\}$; saturating it would incorrectly change
the Fourier coefficient at $(0,1)$.  In the window $[-1,1]^2$,
the two one-frequencies $(1,1)$ and $(-1,1)$ generate the
pure-base witness $(2,0)$, even though neither nonzero pure-base
frequency in that window belongs to their integer span.  Thus
the intersection test must inspect the generated subgroup, not
only the displayed one-frequencies or the base frequencies
inside the window.

\paragraph{A computable finite Fourier approximation.}
Use the product Fej\'er kernel $P_M$ on $\mathcal T$, with
frequency window
$W_M=[-M,M]^{d+r}\cap\Z^{d+r}$, and set
$g_{\theta,M}=P_M*g_\theta$.  Let $L_g$ be a computable common
Lipschitz bound for the scalar functions $g_\theta$ in these
real torus coordinates.  Then
\[
 \|g_\theta-g_{\theta,M}\|_\infty\le e_M,
 \qquad
 e_M=L_g(d+r)
          \frac{1+2\log(M+1)}{4(M+1)}\longrightarrow0.
 \tag{C38g}
\]
To check the constant, on $[-1/2,1/2]$ the one-dimensional
kernel satisfies
\[
 P_M(t)\le\min\left(M+1,\frac1{4(M+1)t^2}\right),
\]
using $|\sin\pi t|\ge2|t|$.  Splitting the first absolute
moment at $1/(2(M+1))$ gives
$\int|t|P_M(t)\,dt\le(1+2\log(M+1))/(4(M+1))$.
The product kernel and the bound of Euclidean distance by the sum
of coordinate distances prove (C38g).

Write
\[
       g_{\theta,M}(x)=
               \sum_{\xi\in W_M}a_{M,\theta}(\xi)e^{2\pi i\xi x}.
\]
For every admissible pattern from (C38f), its Haar average is the
finite real expression
\[
           S_{M,b}(\theta)=
                   \sum_{\xi\in W_M,\ b(\xi)=1}a_{M,\theta}(\xi).
 \tag{C38h}
\]
Uniform interval quadrature computes these coefficients to any
specified precision; certified derivative bounds control the
quadrature error.  Finite nets on the unit sphere and the unit
circles, with the computable parameter Lipschitz bounds, then
give a certified lower bound $m_M$ for all the finitely many
expressions (C38h), uniformly in $\theta$.  The nets and
precision may be chosen so that
\[
 0\le \min_{b,\theta}S_{M,b}(\theta)-m_M\le 2^{-M}.
 \tag{C38i}
\]
For example, bound the parameter Lipschitz constant of each
coefficient by the corresponding constant for $g_\theta$;
the sum has at most $|W_M|$ terms.  Choose a net fine enough
that this sum varies by at most $2^{-M-1}$ on each net cell,
and compute each coefficient to error at most
$2^{-M-1}/|W_M|$.  Taking the smallest interval lower endpoint
gives (C38i), after an immaterial further halving of these
tolerances if needed for both endpoint errors.

Increase $M$ until the computed positive number $m=m_M$ satisfies
\[
                              m>4e_M.
 \tag{C38j}
\]
This search terminates.  Every admissible pattern comes from
some $G_b$, and (C38d)--(C38g) show that all (C38h) are at
least $\Delta-e_M$.  Equations (C38i) and $e_M\to0$ therefore
force (C38j) for sufficiently large $M$.

\paragraph{Making the actual finite graph pattern admissible.}
There is a finite and explicit reduction test; merely checking
the pure-base frequencies lying in $W_M$ would not suffice.
For every subset $S\subset W_M$, form $D_S=\sum_{\xi\in S}\Z\xi$.
If $D_S\cap(\Z^d\times0)\ne0$, compute one nonzero witness
$(\gamma_S,0)$ in this intersection.  Let $\Gamma_M$ be the
finite collection of all these nonzero $\gamma_S$.

Take the actual graph measure $\nu_q$ from C37, with its moving
lattice characters determined by the chosen other-factor samples
and their arbitrary lifts.  Mark a frequency $\xi\in W_M$ by
one exactly when it is trivial on the underlying finite graph
group.  Each such frequency has Fourier coefficient one; every
other frequency has coefficient of absolute value at most
$1/(q-1)$ for the CM rank-one samples.  The subgroup generated
by the one-frequencies is contained in the exact annihilator of
that finite graph group.  Consequently no zero-frequency of the
pattern can belong to it.  If its pure-base intersection were
nonzero, its precomputed witness $\gamma_S\in\Gamma_M$ would
annihilate the first-factor finite torsion group.  Thus the
finite tests
\[
       \gamma|_{T[\mathfrak p_q]}\ne1
                     \quad\text{for every }\gamma\in\Gamma_M
 \tag{C38k}
\]
force the actual pattern to satisfy both lines of (C38f).
These are computable trace-dual residue tests using (C37p).
Each fails at only finitely many primes, since every witness is
nonzero.  One may either test each candidate prime directly or
exclude the primes dividing the cleared nonzero coordinate
norms of these finitely many trace-dual elements.  The latter
also yields a finite numerical bound.  No saturation is used
when generating $D_S$ or extracting the witnesses.

Let
\[
             C_M=|W_M|B_g.
\]
This bounds $\sum_{\xi\in W_M}|a_{M,\theta}(\xi)|$ uniformly.
For a prime passing (C38k), the actual graph average and the
Haar average with its admissible finite pattern differ by at
most $C_M/(q-1)$.  Therefore
\[
        \int g_\theta\,d\nu_q
                  \ge m-e_M-\frac{C_M}{q-1}.
 \tag{C38l}
\]
In particular, the computable additional condition
$q-1\ge4C_M/m$, together with (C38j), makes the right side
at least $m/2$.  This comparison works for the same rational
prime in all the factors.  The moving characters and their
possibly very large lattice coefficients remain exactly those
of the actual samples.

\paragraph{Choosing actual other-factor samples and certifying rank.}
The preceding calculation used the limiting translations $t_i$.
The theta estimates also give an effective uniform continuity
bound when these translations vary in a compact neighborhood.
Choose a computable $\rho>0$ so that replacing each $t_i$ by
$t'_i$ with $\|t'_i-t_i\|<\rho$ changes every function
$g_\theta$ by at most $m/4$ in the supremum norm.  For example,
a computed translation Lipschitz constant gives such a choice
directly.  Let $M_1$ bound the corresponding actual
$|A_\theta|$ and $|A_\theta(0,1)|$ throughout this neighborhood.

Enumerate the allowed primes until all of (C38k),
$q-1\ge4C_M/m$, and the following finite sampling checks hold:
for each $i\ge2$, one of the rank-one sample points has a lift
$t_{i,q}$ within distance $\rho$ of $t_i$.  These latter
checks use finite enumeration with certified coordinates.  More
precisely, compute every candidate's distance to an error less
than $\rho/8$ and accept only a certified strict upper bound
less than $\rho$.  A point within $\rho/2$ is then certainly
accepted.  Each prime receives only this bounded-precision test;
an undecided boundary point is skipped, not subjected to an
unending equality test.  Equivalently, one may dovetail the
primes, their finite candidate tuples, and the approximation
precision.  Equidistribution into the smaller open balls of
radius $\rho/2$ proves eventual acceptance.  The reduction
tests fail only finitely often and $q$ is unbounded, so this
specified prime search terminates as well.

For those chosen samples, (C38l) and the translation error give
\[
                 \int g_{\theta,\mathrm{actual}}\,d\nu_q\ge m/4
                          \quad(\theta\in\Sigma).
 \tag{C38m}
\]
Since
\[
 (|a|^2-|b|^2)^2
       \le4M_1^2\bigl(|a|-|b|\bigr)^2
       \le4M_1^2|\zeta a-b|^2\quad(|\zeta|=1),
\]
the exact first-factor columns in (C37j) satisfy the explicit
finite certificate
\[
             J_qJ_q^*\ \ge
             \frac{m}{16M_1^2}|S_{1,q}|\,\operatorname{Id}.
 \tag{C38n}
\]
Thus their minimum-norm right inverse has coefficient-sum bound
$4M_1m^{-1/2}\|w\|$.  Converting every factor's character
difference back to literal incidence indicators multiplies this
bound by at most the product of the actual bounds
$2(q_i^{a_i}-1)$ from (A15j).  The fixed other-factor differences
are expanded too; they are not being treated as single permitted
incidence indicators.

\paragraph{What a finite certificate contains.}
A successful run records the window $W_M$, the exact integer
normal-form tests and witnesses $\Gamma_M$, the Fourier and
parameter-grid error bounds, the positive certified value $m$,
the constants $e_M,C_M,\rho,M_1$, the selected prime and its
residue checks, and the selected other-factor sample points.
These finite data verify (C38n) by the displayed inequalities.
Compactness is used to prove that the search terminates, not to
replace any numerical verification step in the resulting
certificate.  A simpler implementation may instead dovetail
primes, actual minors in (A15e), and the precision of their
certified theta approximations, accepting only a determinant
interval separated from zero.  It must not refine one possibly
zero minor indefinitely.  C37 likewise proves termination of
that dovetailed search, but (C38n) has the advantage
of retaining every phase and every lift uniformly.

Finally the exact finite-transform modulation identity and the
fixed nonzero baseline remain the hypotheses of A15 and C37.
The certificate proves a first-derivative spanning statement for
that specified family.  It does not on its own identify the
inverse with a cohomological correspondence, change a typed
archive datum, or remove its canonical period and cover-coordinate
normalizations.
